\documentclass[10pt,a4paper]{amsart}
\usepackage[margin=19mm]{geometry}
\usepackage{amsmath,amssymb,mathtools}
\usepackage[scr=rsfs,cal=euler]{mathalfa}
\usepackage{xfrac}
\usepackage[unicode,hidelinks,bookmarks=false]{hyperref}
\newcommand{\ie}{i.e.\ }
\newcommand{\cf}{cf.\ }
\newcommand{\resp}{respectively\ }
\newcommand{\Subsection}{\subsection}
\providecommand{\nobreakdash}{\nobreak\hbox{-}}
\setlength{\emergencystretch}{2em}
\multlinegap0pt
\usepackage{bm}
\usepackage{bbm}
\usepackage{dsfont}
\usepackage{xspace}
\usepackage{cancel}
\usepackage{mathtools}
\usepackage{amsthm}
\makeatletter
\@ifundefined{theorem}{\newtheorem{theorem}{Theorem}}{}
\makeatother
\newcommand{\R}{\mathbb{R}}
\title[A Nonlocal $p$-Laplacian Interface Model with Sharp Interfac]{A Nonlocal $p$-Laplacian Interface Model with Sharp Interface}
\author{Kehan Shi, Zuoqiang Shi, Tangjun Wang}
\date{Source: arXiv:2606.00594v1 (CC BY 4.0)}
\begin{document}
\maketitle

\noindent\emph{Source and licence.} A Nonlocal $p$-Laplacian Interface Model with Sharp Interface. Version arXiv:2606.00594v1 (CC BY 4.0), distributed under the Creative Commons Attribution 4.0 International licence, as recorded in the arXiv licence metadata for this exact version and shown on the abstract page https://arxiv.org/abs/2606.00594v1. Full rights notice: licenses/arxiv-2606.00594-CC-BY-4.0.txt.\par\smallskip

\noindent\textit{Editorial scope.} Counted unit: the Theorem whose source label is \texttt{th:existence\_nonlocal} in the article (source0.tex lines 558--560, source label \texttt{th:existence\_nonlocal}), delivered together with the article's own complete original proof of it (source0.tex lines 562--638, one \texttt{proof} environment of 77 lines). No mathematics is written, repaired or re-derived: statement and proof are the article's own text line for line. Required context retained uncounted: eq:Ku:3 (lines 412--414); eq:poincare (lines 491--499). Every definition, display and label of those lines is retained unchanged. Omitted from this package and not redistributed: 1--411, 415--490, 500--557, 561--561, 639--1038. Nothing inside a retained excerpt is omitted or altered: every retained body is delivered exactly as the article's lines are written, so no omission receipt of any kind accompanies this package. Citations: the retained excerpts carry no citation command at all, so no bibliography entry of the article is transcribed and no reference is redistributed. Reference adaptation: none was needed and none was made; every reference inside the delivered unit resolves inside it, so no unresolved reference or citation is produced. Printed slips kept literal: no printed word, symbol, hypothesis or number of the counted statement or of its proof was altered. Layout and class: the article's own class is \texttt{article} with packages including a4wide, babel, graphicx, float, hyperref, amssymb, empheq, bbold, enumitem, geometry, caption, amsthm, bbm, tikz; the wrapper is the lane's pinned \texttt{amsart} editorial wrapper at 10pt on A4, which restates the theorem-like environments the retained text needs and adds the article's own macro definitions that the retained text needs (\texttt{\textbackslash R}), with extra wrapper packages (bm, bbm, dsfont, xspace, cancel, mathtools). The delivered numbering is the unit's own. Assets: no retained span contains a figure environment, a graphics-inclusion command, a PDF-page inclusion, a table or a TikZ picture; no image file is copied into this package, the package declares no asset dependency and no image file is redistributed with it. Licence: version arXiv:2606.00594v1 of this article is distributed under the Creative Commons Attribution 4.0 International licence, as recorded in the arXiv licence metadata for this exact version and shown on the abstract page https://arxiv.org/abs/2606.00594v1. A full rights notice is delivered at licenses/arxiv-2606.00594-CC-BY-4.0.txt. Characters: every retained excerpt is pure ASCII, so the delivered engine, XeLaTeX with its default Latin Modern text font, reports no missing character.
  \begin{equation}\label{eq:Ku:3}
    K_{\delta,i}\ast u_n\rightarrow  K_{\delta,i}\ast u,\quad \mbox{in } L^p(\Gamma),
  \end{equation}

  \begin{align}\label{eq:poincare}
  \begin{split}
    \frac{1}{C}\int_{\Omega} |u|^pdx
    \leq \int_{\Omega_1}\int_{\Omega_1}R_\delta(x-y)|u(y)-u(x)|^pdydx
    &+\int_{\Omega_2}\int_{\Omega_2}R_\delta(x-y)|u(y)-u(x)|^pdydx\\
    &+\int_\Gamma\left|K_{\delta,1}\ast u-K_{\delta,2}\ast u\right|^pdS
    +\left|\int_{\Omega} udx\right|^p.
    \end{split}
  \end{align}

\begin{theorem}\label{th:existence_nonlocal}
  Let $\lambda_1, \lambda_2>0$, $f\in L^2(\Omega)$, and $g\in L^2(\Gamma)$. For any $\delta>0$, the nonlocal functional $\mathcal{E}_\delta$ admits a unique minimizer in $L^2_0(\Omega)$.
\end{theorem}

\begin{proof}
  Let $u\in L^2_0(\Omega)$.
  The definition of $\mathcal{E}_\delta$ and nonlocal Poincar\'{e}'s inequality \eqref{eq:poincare} imply
  \begin{align*}
    \int_\Omega |u|^2dx
    \leq C\mathcal{E}_\delta(u)
    +C\int_{\Omega}fudx + C\int_\Gamma g K_{\delta,1}\ast u dS.
  \end{align*}
  Utilizing Cauchy's inequality and the fact
  $\int_\Gamma |K_{\delta,1}\ast u|^2dS\leq C\int_{\Omega_1}|u|^2dx$,
  we have
  \begin{align*}
  \int_{\Omega}fudx + \int_\Gamma gK_{\delta,1}\ast u dS
  \leq& \frac{4}{\epsilon}\int_\Omega |f|^2dx
  +\epsilon\int_\Omega |u|^2dx
  +\frac{4}{\epsilon}\int_\Gamma |g|^2dS
  +C\epsilon \int_{\Omega_1}|u|^2dx.
  \end{align*}
  Combining two inequalities and choosing a small $\epsilon$ yield the coercivity of $\mathcal{E}_\delta$, i.e.,
    \begin{align}\label{eq:coercivity}
  \begin{split}
    \int_\Omega|u|^2dx
    \leq C\mathcal{E}_\delta(u)+C\int_\Omega |f|^2dx + C\int_\Gamma |g|^2dS.
  \end{split}
  \end{align}
There exists a minimizing sequence $u_k\in L_0^2(\Omega)$, such that
\[
\lim_{k\rightarrow\infty}\mathcal{E}_\delta(u_k)=\inf_{L^2(\Omega)} \mathcal{E}_\delta(u),
\]
and $\mathcal{E}_\delta(u_k)$ is uniformly bounded in $L^2(\Omega)$.
\eqref{eq:coercivity} further implies the existence of a subsequence of $\{u_k\}$ (still denoted by itself) and a function $u\in L^2(\Omega)$, such that
\begin{align*}
  u_{k} \rightharpoonup {u}, \quad \mbox{in } L^2(\Omega).
\end{align*}
By \eqref{eq:Ku:3},
  \[
  K_{\delta,i}\ast u_k\rightarrow K_{\delta,i}\ast u, \quad \mbox{in } L^2(\Gamma),
  \]
  as $k\rightarrow \infty$.
 By the weak lower semi-continuity,
\[
\mathcal{E}_\delta(u)
\leq \liminf_{k\rightarrow\infty}\mathcal{E}_\delta(u_k).
\]
Consequently,
\[
\inf_{L^2(\Omega)} \mathcal{E}_\delta(u)
\leq \mathcal{E}_\delta(u)
\leq \liminf_{k\rightarrow\infty}\mathcal{E}_\delta(u_k)
=\inf_{L^2(\Omega)} \mathcal{E}_\delta(u).
\]
Thus $u$ is a minimizer of $\mathcal{E}_\delta$ in $L^2(\Omega)$.
Clearly, $u\in L^2_0(\Omega)$.

To prove the uniqueness of the minimizer, we only need to show that
\[
E_\delta(u)=
\frac{1}{\delta^2}\sum_{i=1,2}\int_{\Omega_i}\int_{\Omega_i}\Lambda(x,y) R_\delta(x-y)|u(y)-u(x)|^2dydx
+\frac{1}{\delta}
\int_\Gamma\left|K_{\delta,1}\ast u-K_{\delta,2}\ast u\right|^2dS
\]
is strictly convex on $L^2_0(\Omega)$.
In fact, for any $u,v\in L^2_0(\Omega)$, $w:=u-v$, and $0<t<1$,
it follows from the elementary equality
\[
t|x|^2+(1-t)|y|^2-|tx+(1-t)y|^2=t(1-t)|x-y|^2,\quad x,y\in\R,
\]
and nonlocal Poincar\'{e}'s inequality \eqref{eq:poincare}
that
\begin{align*}
tE_\delta(u)+(1-t)E_\delta(v)-E_\delta(tu+(1-t)v)
=t(1-t)E_\delta(w)
\geq Ct(1-t)\int_\Omega |w|^2dx>0,
\end{align*}
if $u\neq v$.
This completes the proof.
\end{proof}

\end{document}
